\section{Ejecución del bloque IV: P13--P16}
\label{sec:bloque4}
\textbf{Fecha: 18 de octubre de 2026.} P13--P15 están implementados y probados
localmente. P16 incorpora instrumentación y un perfil del runtime con catálogo
e índices reales y modelos de control; sus mediciones con modelos/proveedor
reales y memoria completa del proceso siguen pendientes. P04/P10/P12 y E13
no se consideran cerrados por avanzar en este bloque.

\subsection{P13: puntuación explicable y pesos configurables}
\textbf{Antes.} Se combinaban scores de workers normalizados por mínimo/máximo
con una bonificación RRF independiente de los pesos. El resultado se multiplicaba
por cien y podía superar cien, sin ser una probabilidad. Una rama con peso cero
podía seguir aportando por la bonificación.

\textbf{Ahora.} \texttt{scoring.py} define una fusión por posición:
\[
R_m(d)=\begin{cases}\frac{61}{60+r_m(d)},&d\text{ aparece en la rama }m,\\
0,&\text{en otro caso},\end{cases}
\qquad S(d)=w_L R_L(d)+w_A R_A(d).
\]
La posición comienza en uno. Los pesos efectivos suman uno entre ramas con peso
positivo que aportan resultados. Se exponen rango, score bruto del worker,
score por posición, peso y contribución por canción; las contribuciones suman
el score final. Los empates se resuelven por ID, de forma determinista.

La API versión 2 elimina \texttt{score\_pct} y añade el método
\texttt{weighted\_normalized\_rrf\_v2}, rango $[0,1]$ e indicador de que el score
no es probabilidad. La interfaz presenta decimales y explica el significado.
Las decisiones del juez pueden priorizar una recomendación frente a otra con
mayor score: se declara esa política y su motivo por resultado.

Los pesos opcionales de solicitud son dos números finitos entre cero y uno que
suman uno. Se distinguen pesos solicitados y efectivos, con ajuste explícito
ante ramas ausentes, vacías o fallidas. Peso cero desactiva la recuperación de
esa rama, sin eliminar conceptos o restricciones que debe comprobar el juez.
No se inventa una modalidad cuando el plan carece de consultas para ella.

\textbf{Ventaja y contrapartida.} Se elimina el porcentaje engañoso y el aporte
de ramas con peso cero. La nueva fusión evita comparar magnitudes de espacios
distintos, pero descarta diferencias de magnitud dentro de cada ranking y puede
cambiar el orden respecto a la versión anterior. Los primeros puestos tienden
a valores cercanos a uno; ello no demuestra relevancia alta. Scores de versiones,
consultas y ventanas distintas no se comparan como probabilidades. La superioridad
musical de esta fórmula sigue pendiente de P04/P12; no se infiere de sus contratos.

\subsection{P14: filtros duros por metadatos}
\textbf{Antes.} No había selección exacta por artista, género o fechas.
\textbf{Ahora.} \texttt{search\_options.py} valida artista y género completos,
ignorando mayúsculas y espacios exteriores y normalizando Unicode NFC, sin
eliminar acentos ni hacer coincidencias parciales. El género procede de
\texttt{tag}; no se infiere del audio. Los años son enteros y los límites son
inclusivos. Los filtros se combinan con AND; años ausentes o inválidos no
satisfacen un filtro temporal. Columnas opcionales ausentes no rompen los
controles de catálogo.

\texttt{GET /catalog/options} ofrece nombres y límites de años sin letras ni
llamadas a modelos. Con filtros activos se recupera hasta el tamaño del catálogo
antes de filtrar cada rama y fusionar la ventana elegible. Así, una canción en
posición 25 del ranking global puede ser la primera coincidencia filtrada.
El juez solo recibe candidatos elegibles. Cero coincidencias de metadatos se
resuelve antes del router, sin consumir modelos, y queda distinguido de una
consulta semántica sin resultados utilizables.

\textbf{Evidencia y límites.} Pruebas HTTP cubren esa canción fuera del top-10,
intersección, tipos inválidos, límites inclusivos y ausencia de coincidencias.
La operación usa un catálogo validado y estable durante el proceso; P17 sigue
siendo necesario para actualización coordinada. Recuperar todos los candidatos
es sencillo para 258 canciones, pero aumenta trabajo y memoria al crecer el
catálogo. Evaluar preselección nativa en el índice o índices de metadatos dentro
de P16/P17 antes de escalar. La completitud de metadatos tampoco certifica su
exactitud ni su género acústico.

\subsection{P15: evidencia visible y alcance por candidato}
\textbf{Antes.} La respuesta descartaba letras completas y la tarjeta mostraba
principalmente título, score y un estado, sin fragmentos ni contribuciones.
\textbf{Ahora.} \texttt{result\_evidence.py} selecciona hasta 320 caracteres.
Prioriza citas estructuradas que ocurren literalmente en la letra. Sin esas
citas, usa una heurística léxica de las consultas positivas o un fragmento
inicial, rotulados como orientativos. Conserva offsets del texto original y
advierte recortes. Las solicitudes puramente acústicas no reciben una supuesta
prueba lírica. No se envía la letra completa como parte de cada resultado.

Cada tarjeta muestra contribuciones, metadatos, alcance de letra/audio, motivo
del juez cuando existe y distinción entre recomendación evaluada y alternativa
recuperada. Se conserva el estado de error de un candidato realmente evaluado;
no se reemplaza por un falso estado de candidato nunca examinado. Se usa solo
evidencia de la versión de plan correspondiente. Las decisiones heredadas sin
citas mantienen esa limitación explícita. El contenido variable se inserta como
texto DOM, incluyendo títulos, citas y mensajes.

\textbf{Evidencia y límites.} Pruebas cubren citas reales/inventadas, offsets
Unicode, ausencia de aprobación heredada en alternativas y contexto acústico.
La web se verificó en un navegador con fixture visiblemente marcado como
control: catálogo de 258 canciones, rankings y juez simulados, filtros
combinados, pesos personalizados, año inválido y cero coincidencias. La captura
\texttt{validation/block4/ui-desktop.jpg} permite revisar el resultado. Hay CSS
adaptativo; no se certifica una revisión visual móvil. Una coincidencia textual
o un fragmento orientativo no prueban cumplimiento semántico: E13/P10 siguen abiertos.

\subsection{P16: instrumentación y decisiones basadas en mediciones}
\textbf{Antes.} No había desglose por construcción del grafo, ejecución,
presentación, codificación, índice y contención. \textbf{Ahora.} La API informa
tiempos del servicio y nodos. Las rondas de recuperación registran caché de la
solicitud, perfil de ejecución de workers y fallos. En un acierto de caché no se
atribuyen otra vez los tiempos del worker anterior. Los workers miden acceso a
recursos, espera del bloqueo, codificación/caché, búsqueda FAISS y RRF interno.
Los perfiles usan contexto de ejecución separado para solicitudes concurrentes.
Los tiempos internos no incluyen necesariamente serialización HTTP ni red.

\texttt{evaluation/profile.py} ejecuta una primera consulta, 20 consultas
calientes secuenciales y 20 con dos y con cuatro hilos. Usa los índices FAISS
y metadatos reales; los encoders reconstruyen vectores existentes y router/juez
son controles. Por tanto, se mide el runtime local, no inferencia Nomic/CLAP ni
servicio Groq. Se fija un hilo OpenMP por tarea solo en el experimento. Reporta
medias, mediana, p95, dispersión, tasa de solicitudes, contención y cachés.
Cinco solicitudes adicionales miden asignaciones Python con
\texttt{tracemalloc} en una fase separada, para no sesgar los tiempos calientes
frente a la primera consulta sin trazado. El total es 66 solicitudes de control.

El experimento verificó identidad de resultados entre las condiciones y
reutilización de cachés. La memoria trazada excluye tensores nativos, FAISS y
GPU, y no equivale a RSS; se informa separadamente el payload de vectores del
catálogo. Las métricas de memoria completa y modelos reales permanecen nulas.
Los valores y hashes reproducibles están en \texttt{validation/block4/profile.json}.

\textbf{Decisión y pendiente.} Se mantienen grafos por solicitud y cachés
existentes acotadas; no se introduce un grafo compartido, otra caché global ni
batching por una aceleración aparente en modelos de control. Un grafo reutilizado
exigiría separar su caché de cada solicitud y medir invalidación/contención.
Comparar con máquina de estados simple solo si la compilación pesa realmente en
la carga final. P16 continúa abierto para carga con modelos, memoria completa y
optimización demostrada; E04/E09 y P17/P18 conservan su trazabilidad.

\subsection{Archivos, compatibilidad y verificación}
\texttt{app.py} integra score, filtros, evidencia y mediciones; \texttt{cli.py}
admite peso lírico, artista, género y años. Se añaden
\texttt{scoring.py}, \texttt{search\_options.py}, \texttt{result\_evidence.py} y
\texttt{profiling.py}. Se instrumentan \texttt{workflow.py} y ambos workers.
La plantilla contiene controles y tarjetas actualizadas. Las pruebas nuevas
están en \texttt{tests/test\_block4.py}; los informes y capturas, en
\texttt{validation/block4/}. \texttt{scripts/preview\_block4.py} es exclusivamente
un fixture local de interfaz, señalado como simulación; no es configuración de
producción ni evaluación musical.

Se ejecutaron 115 pruebas en CP1252 y 115 en UTF-8. Se mantienen capturas de
36 simulaciones sin refinamiento y 36 con él, sin modificar las consultas
históricas ni sustituir juicios humanos. La evaluación real sigue bloqueada
cuando falta la clave. El perfil adicional no descarga pesos ni llama al
proveedor. No hay dependencias nuevas. Las fuentes TeX se verifican por
inclusiones y referencias; no se certifica una nueva compilación PDF.

Los clientes deben migrar de \texttt{score\_pct} a \texttt{score} y sus
contribuciones; el endpoint sigue admitiendo solicitudes antiguas con solo
\texttt{query}. La firma de consola comparte los mismos contratos de opciones.
Los informes anteriores conservan fórmulas y snapshots de su fecha: no comparar
scores numéricos entre métodos como si compartieran escala semántica.

El siguiente bloque de implementación es P17--P20. Permanecen pendientes las
validaciones reales de P04/P10/P12/P16 y la fidelidad E13. El bloque IV aporta
transparencia y controles, sin atribuirse relevancia o rendimiento de modelos
todavía no medidos.
